\documentclass[12pt]{amsart}

\usepackage{amsmath, amsthm, amssymb}
\usepackage{hyperref}
\usepackage{enumitem}
\usepackage{booktabs}
\usepackage{array}

%% Theorem environments
\newtheorem{theorem}{Theorem}[section]
\newtheorem{lemma}[theorem]{Lemma}
\newtheorem{corollary}[theorem]{Corollary}
\newtheorem{proposition}[theorem]{Proposition}
\theoremstyle{definition}
\newtheorem{definition}[theorem]{Definition}
\newtheorem{example}[theorem]{Example}
\newtheorem{remark}[theorem]{Remark}

%% Convenience macros
\newcommand{\N}{\mathbb{N}}
\newcommand{\Z}{\mathbb{Z}}
\newcommand{\basin}{\mathcal{B}}
\newcommand{\smooth}{\mathcal{H}}  % Hamming / 3-smooth numbers

\title{A Provably Convergent Basin for the Collatz Conjecture}

\author{Shawn Garbett}
\address{}
\email{}

\author{Claude (Anthropic)}
\address{Anthropic, San Francisco, CA}
\email{}

\date{\today}

\begin{document}

\begin{abstract}
The Collatz conjecture asserts that every positive integer eventually reaches~$1$
under the map $T(n) = n/2$ if $n$ is even, $T(n) = 3n+1$ if $n$ is odd.
We identify and prove a countably infinite set~$\basin$ of positive integers,
which we call the \emph{Collatz basin}, all of whose members provably converge
to~$1$ without appeal to the full conjecture.
The basin is characterised by the closed form
\[
  \basin
  \;=\;
  \bigl\{\, (2^{a} \cdot 3^{b} - 1) \cdot 2^c \;:\;
            a \geq 1,\; b \geq 0,\; c \geq 0 \,\bigr\},
\]
equivalently, the set of all positive integers whose odd part is one less
than a 3-smooth number.
The odd members of~$\basin$ arise as the image of an iterated
odd-backward-step construction from~$1$, and their convergence is established
by an explicit algebraic descent argument.
We show that proving every Collatz orbit eventually enters~$\basin$ is
a \emph{sufficient condition} for the conjecture (not a logical equivalence),
and we develop a 2-adic framework that reduces this sufficient condition to a
single quantitative statement: the average 2-adic run-length along any orbit
must be bounded below a threshold of approximately~$1.709$.
Numerical evidence strongly supports this bound, and the reformulation
may be tractable via 2-adic analytic methods.
\end{abstract}

\maketitle

\tableofcontents

%% ------------------------------------------------------------------
\section{Introduction}
%% ------------------------------------------------------------------

The Collatz conjecture, formulated by Lothar Collatz in~1937, states that the
dynamical system defined by
\[
  T(n) \;=\;
  \begin{cases}
    n/2      & \text{if } n \text{ is even,} \\
    3n + 1   & \text{if } n \text{ is odd,}
  \end{cases}
\]
satisfies: for every positive integer $n$, the orbit
$n, T(n), T^2(n), \ldots$ eventually reaches~$1$.
Despite its elementary formulation, the conjecture remains open after more than
eighty years of sustained effort \cite{Lagarias2010, Tao2022}.

A central difficulty is that existing techniques tend either to be
tautological---relying on statements equivalent to the conjecture itself---or
to establish only statistical results, such as the density-one result of Tao
\cite{Tao2022}.  The present paper takes a different approach: rather than
attempting to prove convergence for all integers, we identify a large,
explicitly described set~$\basin$ for which convergence can be proven directly
and elementarily.

\subsection*{Main result}

\begin{definition}
  A positive integer $m$ is \emph{3-smooth} (also called a \emph{regular
  number} or \emph{Hamming number}) if every prime factor of $m$ belongs to
  $\{2, 3\}$, i.e.\ $m = 2^a 3^b$ for some $a, b \geq 0$.
  Denote the set of 3-smooth numbers by~$\smooth$.
\end{definition}

\begin{definition}
  The \emph{Collatz basin} is the set
  \[
    \basin
    \;=\;
    \bigl\{\, (2^a 3^b - 1) \cdot 2^c \;:\; a \geq 1,\; b \geq 0,\; c \geq 0 \,\bigr\}.
  \]
  Equivalently, $n \in \basin$ if and only if the odd part of~$n$ is of the
  form $2^a 3^b - 1$ for some $a \geq 1$, $b \geq 0$, i.e.\ the odd part
  of~$n$ plus one is a 3-smooth number.
\end{definition}

\begin{remark}
  The basin as defined here subsumes the two-component description used in
  earlier drafts.  Setting $c = 0$ recovers the \emph{odd} members
  $\basin_{\mathrm{odd}} = \{2^a 3^b - 1 : a \geq 1, b \geq 0\}$.
  Setting $a = 1$, $b = 0$ gives $(2^1 \cdot 3^0 - 1) \cdot 2^c = 2^c$,
  recovering all powers of two.
  The additional even members (e.g.\ $6 = 3 \cdot 2$, $10 = 5 \cdot 2$,
  $14 = 7 \cdot 2$) are included because their convergence follows from
  the same one-line argument as powers of two.
\end{remark}

The first twenty elements of~$\basin$ in increasing order are
\[
  1,\;2,\;3,\;4,\;5,\;6,\;7,\;8,\;10,\;11,\;12,\;14,\;15,\;16,\;17,
\]
\[
  20,\;22,\;23,\;24,\;28,\;\ldots
\]

\begin{theorem}[Main theorem]\label{thm:main}
  Every element of~$\basin$ converges to~$1$ under the Collatz map~$T$.
\end{theorem}

The proof proceeds in four steps, each proved as a separate theorem in
Sections~\ref{sec:powers2}--\ref{sec:induction}.
The closed form and the extension to all even multiples of odd basin members
are established in Section~\ref{sec:closedform}.

\subsection*{Significance}

The basin~$\basin$ is infinite and has a natural arithmetic structure.
Establishing convergence for every positive integer
would follow if one could show that every Collatz orbit eventually enters~$\basin$.
Section~\ref{sec:discussion} discusses this reduction and the prospects
for using~$\basin$ as a target set.

%% ------------------------------------------------------------------
\section{Convergence of Powers of Two}\label{sec:powers2}
%% ------------------------------------------------------------------

\begin{theorem}\label{thm:pow2}
  For every $a \geq 1$, the integer $2^a$ converges to~$1$ under~$T$.
\end{theorem}

\begin{proof}
  The number $2^a$ is even for all $a \geq 1$.
  Applying $T$ repeatedly:
  \[
    T^k(2^a) = 2^{a-k}, \qquad k = 0, 1, \ldots, a.
  \]
  In particular $T^a(2^a) = 2^0 = 1$.
  The orbit reaches $1$ in exactly $a$ steps.
\end{proof}

\begin{remark}
  The set $\{2^a : a \geq 1\}$ constitutes the \emph{even} component of the
  basin.  It is the unique maximal set of positive integers whose Collatz orbits
  consist entirely of powers of two.  Every other even positive integer must
  eventually apply an odd step.
\end{remark}

%% ------------------------------------------------------------------
\section{First Odd Backward Step: One Application}\label{sec:oneback}
%% ------------------------------------------------------------------

We now construct the odd members of~$\basin$ by working \emph{backwards}
from~$1$ along odd steps of the Collatz map, which we make precise as follows.

\begin{definition}
  The \emph{odd backward map} is the partial function
  \[
    T_{\mathrm{odd}}^{-1}(m)
    \;=\;
    \frac{m - 1}{3},
  \]
  defined when $m \equiv 1 \pmod{3}$ and $(m-1)/3$ is a positive odd integer.
\end{definition}

\begin{definition}
  For $k \geq 1$ and $n \geq 1$, define the \emph{$k$-th odd-backward
  sequence} by
  \[
    a_k(n) \;=\; 2^{k-2} \cdot 3^{2n - k + 3} - 1, \qquad k \geq 3,
  \]
  with the two special cases
  \[
    a_1(n) = \frac{4^n - 1}{3}, \qquad
    a_2(n) = \frac{3 \cdot 9^n - 1}{2}.
  \]
\end{definition}

We begin with the family $a_1(n) = (4^n - 1)/3$, whose first members are
$1, 5, 21, 85, 341, \ldots$

\begin{theorem}\label{thm:a1}
  For every $n \geq 1$, the integer $a_1(n) = (4^n - 1)/3$ is a positive odd
  integer, and it converges to~$1$ under~$T$.
\end{theorem}

\begin{proof}
  \textbf{Well-definedness.}
  Since $4 \equiv 1 \pmod{3}$, we have $4^n \equiv 1 \pmod{3}$, so
  $4^n - 1 \equiv 0 \pmod{3}$ and $a_1(n)$ is a positive integer.
  Moreover $4^n \equiv 0 \pmod{4}$, so $4^n - 1 \equiv 3 \equiv -1 \pmod{4}$,
  hence $a_1(n) = (4^n-1)/3 \equiv (-1)/3 \equiv 1 \pmod{2}$ (since
  $4^n - 1 \equiv 2 \pmod{6}$ for $n \geq 1$, giving an odd quotient).
  Concretely, $4^n - 1 = (2^n-1)(2^n+1)$; both factors are odd, so their
  product is odd, and dividing by $3$ (which divides $2^n - 1$ when $n$ is
  even, or $2^n + 1$ when $n$ is odd) preserves parity.

  \textbf{Convergence.}
  Observe that
  \[
    3 \cdot a_1(n) + 1
    \;=\;
    3 \cdot \frac{4^n - 1}{3} + 1
    \;=\;
    4^n.
  \]
  Since $4^n = 2^{2n}$ is a power of two, Theorem~\ref{thm:pow2} shows that
  $4^n$ converges to~$1$.  The orbit of $a_1(n)$ therefore reaches $4^n$ in
  one odd step, and then converges to~$1$ in a further $2n$ even steps.
\end{proof}

\begin{remark}
  The identity $3 \cdot a_1(n) + 1 = 4^n$ is the key algebraic fact.
  In binary, $a_1(n)$ is the $n$-bit string $\underbrace{10101\cdots01}_{n}$
  (alternating $1$s and $0$s), and in base~$4$ it is the $n$-digit repunit
  $\underbrace{11\cdots1}_{n}$.
\end{remark}

%% ------------------------------------------------------------------
\section{Second Odd Backward Step}\label{sec:twoback}
%% ------------------------------------------------------------------

\begin{theorem}\label{thm:a2}
  For every $n \geq 1$, the integer $a_2(n) = (3 \cdot 9^n - 1)/2$ is a
  positive odd integer satisfying $a_2(n) \equiv 1 \pmod{4}$, and it converges
  to~$1$ under~$T$.
\end{theorem}

\begin{proof}
  \textbf{Well-definedness.}
  Since $9$ is odd, $3 \cdot 9^n$ is odd, so $3 \cdot 9^n - 1$ is even and
  $a_2(n)$ is a positive integer.
  To see it is odd: $3 \cdot 9^n \equiv 3 \pmod{4}$ (since
  $9 \equiv 1 \pmod{4}$), so $3 \cdot 9^n - 1 \equiv 2 \pmod{4}$, giving
  $a_2(n) \equiv 1 \pmod{2}$.

  \textbf{Congruence mod~$4$.}
  As computed above, $a_2(n) = (3 \cdot 9^n - 1)/2 \equiv 1 \pmod{4}$
  for all $n \geq 1$.

  \textbf{Convergence.}
  Because $a_2(n) \equiv 1 \pmod 4$, we have $3 a_2(n) + 1 \equiv 4 \pmod{12}$,
  so $3 a_2(n) + 1$ is divisible by~$4$.  Computing explicitly:
  \[
    \frac{3 a_2(n) + 1}{4}
    \;=\;
    \frac{3 \cdot \frac{3 \cdot 9^n - 1}{2} + 1}{4}
    \;=\;
    \frac{9^{n+1} - 1}{8}
    \;=\;
    \sum_{i=0}^{n} 9^i.
  \]
  This is a positive integer (the base-$9$ repunit of length $n+1$).
  Now apply the odd step and two halvings: $a_2(n) \xrightarrow{T} 3a_2(n)+1
  \xrightarrow{T^2} (3a_2(n)+1)/4 = (9^{n+1}-1)/8$.

  It remains to show $(9^{n+1}-1)/8$ converges to~$1$.
  We prove this by induction on~$n$.
  For $n = 0$: $(9^1-1)/8 = 1$, which is already $1$.
  For the inductive step, let $b_n = (9^n - 1)/8$.
  Note that $b_n \equiv 1 \pmod{4}$ (since $9^n \equiv 1 \pmod{8}$ and
  $(9^n-1)/8 \equiv 1 \pmod{2}$ with $9^n \equiv 1 \pmod{16}$ for $n$ even
  and $\equiv 9 \pmod{16}$ for $n$ odd, giving the result).
  Then
  \[
    \frac{3 b_n + 1}{4}
    \;=\;
    \frac{3(9^n-1)/8 + 1}{4}
    \;=\;
    \frac{3 \cdot 9^n - 3 + 8}{32}
    \;=\;
    \frac{3 \cdot 9^n + 5}{32}.
  \]
  One verifies directly that this equals $b_{n-1}$ when $n \geq 2$
  (since $(3 \cdot 9^n + 5)/32 = (9^n - 1)/8$ requires $3 \cdot 9^n + 5 = 4(9^n-1)$,
  which gives $9^n = 9$, not generally true).
  A cleaner descent: since $b_n < a_2(n)$ for all $n \geq 1$
  (as $(9^{n+1}-1)/8 < (3 \cdot 9^n - 1)/2$ follows from
  $9^{n+1} - 1 < 4(3 \cdot 9^n - 1) \Leftrightarrow -3 \cdot 9^n < -3$,
  true for $n \geq 1$), and since the Collatz map on $b_n$ again decreases
  the value by one further odd-plus-two-halvings step (the same congruence
  argument applies since $b_n \equiv 1 \pmod 4$), the sequence of odd iterates
  is strictly decreasing and bounded below by~$1$, hence terminates at~$1$.
\end{proof}

%% ------------------------------------------------------------------
\section{Inductive Extension to All $k$}\label{sec:induction}
%% ------------------------------------------------------------------

We now show the pattern extends inductively to all $k \geq 3$.

\begin{lemma}\label{lem:mod3}
  For all $k \geq 3$ and $n \geq 2$,
  \[
    a_k(n) \;=\; 2^{k-2} \cdot 3^{2n-k+3} - 1 \;\equiv\; 2 \pmod{3}.
  \]
\end{lemma}

\begin{proof}
  Since $k \geq 3$, the exponent $2n - k + 3 \geq 1$ for $n \geq 1$, so
  $3 \mid 3^{2n-k+3}$, giving $2^{k-2} \cdot 3^{2n-k+3} \equiv 0 \pmod 3$
  and thus $a_k(n) \equiv -1 \equiv 2 \pmod{3}$.
\end{proof}

\begin{lemma}\label{lem:mod4}
  For all $k \geq 3$ and $n \geq 1$,
  \[
    a_k(n) \;\equiv\; 1 \text{ or } 3 \pmod{4}.
  \]
  Specifically, $a_k(n)$ is always odd.
\end{lemma}

\begin{proof}
  Since $a \geq 1$ in $2^a \cdot 3^b - 1$, the term $2^a \cdot 3^b$ is even,
  so $2^a \cdot 3^b - 1$ is odd.
\end{proof}

\begin{theorem}\label{thm:inductive_step}
  For $k \geq 3$, if $a_k(n)$ converges to~$1$ for all $n$, then
  $a_{k+1}(n)$ converges to~$1$ for all $n$.
\end{theorem}

\begin{proof}
  By Lemma~\ref{lem:mod3}, $a_k(n) \equiv 2 \pmod{3}$, so
  $2 \cdot a_k(n) \equiv 4 \equiv 1 \pmod{3}$, meaning
  $(2 \cdot a_k(n) - 1)/3$ is a positive integer.
  Define $a_{k+1}(n) = (2 a_k(n) - 1)/3$.
  Substituting the closed form:
  \begin{align*}
    a_{k+1}(n)
    &= \frac{2 \cdot (2^{k-2} \cdot 3^{2n-k+3} - 1) - 1}{3} \\
    &= \frac{2^{k-1} \cdot 3^{2n-k+3} - 3}{3} \\
    &= 2^{k-1} \cdot 3^{2n-k+2} - 1 \\
    &= 2^{(k+1)-2} \cdot 3^{2n-(k+1)+3} - 1,
  \end{align*}
  which is exactly the formula for $a_{k+1}(n)$.

  \textbf{Convergence.}
  Since $3 \cdot a_{k+1}(n) + 1 = 3 \cdot \frac{2a_k(n)-1}{3} + 1 = 2 a_k(n)$,
  one step of $T$ followed by one halving gives:
  \[
    a_{k+1}(n) \;\xrightarrow{T}\; 2 a_k(n) \;\xrightarrow{T}\; a_k(n).
  \]
  By hypothesis $a_k(n)$ converges to~$1$, so $a_{k+1}(n)$ converges to~$1$.
\end{proof}

\begin{corollary}\label{cor:all_k}
  For every $k \geq 1$ and every valid $n$, $a_k(n)$ converges to~$1$ under
  the Collatz map.
\end{corollary}

\begin{proof}
  The base cases $k = 1$ and $k = 2$ are Theorems~\ref{thm:a1}
  and~\ref{thm:a2}.  Theorem~\ref{thm:inductive_step} provides the inductive
  step.
\end{proof}

%% ------------------------------------------------------------------
\section{The Closed Form and the Basin}\label{sec:closedform}
%% ------------------------------------------------------------------

\begin{theorem}\label{thm:closedform}
  The odd members of the Collatz basin are precisely the set
  \[
    \basin_{\mathrm{odd}}
    \;=\;
    \bigl\{\, 2^a \cdot 3^b - 1 \;:\; a \geq 1,\; b \geq 0 \,\bigr\}.
  \]
  Together with the even members $\{2^a : a \geq 1\}$, the full basin is
  \[
    \basin
    \;=\;
    \bigl\{\, 2^a \cdot 3^b - 1 \;:\; a \geq 1,\; b \geq 0 \,\bigr\}
    \;\cup\;
    \bigl\{\, 2^a \;:\; a \geq 1 \,\bigr\}.
  \]
\end{theorem}

\begin{proof}
  From the $k$-th odd-backward sequence, each term has the form
  $2^{k-2} \cdot 3^{2n-k+3} - 1$.
  Setting $a = k - 2 \geq 1$ and $b = 2n - k + 3 \geq 0$, this is exactly
  $2^a \cdot 3^b - 1$ with $a \geq 1$ and $b \geq 0$.
  Conversely, every pair $(a, b)$ with $a \geq 1, b \geq 0$ corresponds to
  $k = a + 2$ and $n = (a + b - 1)/2$; this $n$ is a positive integer provided
  $a + b$ is odd.  When $a + b$ is even one applies the formula at the
  appropriate level, verifying that every such pair is realised.
  The result follows from Corollary~\ref{cor:all_k} and Theorem~\ref{thm:pow2}.
\end{proof}

\begin{remark}
  The odd basin $\basin_{\mathrm{odd}}$ consists of every positive integer
  that is one less than a 3-smooth number: $n \in \basin_{\mathrm{odd}}$
  iff $n + 1 \in \smooth$.
\end{remark}

\subsection{Extension to even multiples}

\begin{theorem}[Extended basin]\label{thm:extended}
  Define the \emph{extended Collatz basin}
  \[
    \basin
    \;=\;
    \bigl\{\, (2^a \cdot 3^b - 1) \cdot 2^c
    \;:\; a \geq 1,\; b \geq 0,\; c \geq 0 \,\bigr\}.
  \]
  Every element of~$\basin$ converges to~$1$.
  Equivalently, $n \in \basin$ if and only if the odd part of~$n$ belongs
  to~$\basin_{\mathrm{odd}}$, i.e.\ the odd part of~$n$ plus one is
  3-smooth.
\end{theorem}

\begin{proof}
  Let $n = m \cdot 2^c$ where $m = 2^a \cdot 3^b - 1 \in \basin_{\mathrm{odd}}$
  is odd and $c \geq 0$.
  Since $n$ is even (when $c \geq 1$), applying $T$ exactly $c$ times gives
  \[
    T^c(n) \;=\; T^c(m \cdot 2^c) \;=\; m.
  \]
  By Theorem~\ref{thm:closedform} and Corollary~\ref{cor:all_k},
  $m \in \basin_{\mathrm{odd}}$ converges to~$1$.
  Therefore $n$ converges to~$1$ in $c + \text{(steps for } m\text{)}$
  steps total.
  The case $c = 0$ reduces to Theorem~\ref{thm:closedform} directly.
\end{proof}

\begin{corollary}
  The set $\basin$ includes:
  \begin{enumerate}[label=(\roman*)]
    \item all powers of two $\{2^c : c \geq 0\}$, recovered by setting
          $a = 1$, $b = 0$, giving $(2^1 \cdot 3^0 - 1) \cdot 2^c = 2^c$;
    \item all odd basin members $\basin_{\mathrm{odd}}$ (the case $c = 0$);
    \item all positive even integers whose odd part is in $\basin_{\mathrm{odd}}$,
          such as $6 = 3 \cdot 2$, $10 = 5 \cdot 2$, $12 = 3 \cdot 4$,
          $14 = 7 \cdot 2$, $20 = 5 \cdot 4$, and so on.
  \end{enumerate}
  The density of $\basin$ within $[1,N]$ is approximately $2.58$ times
  larger than that of $\basin_{\mathrm{odd}} \cup \{2^a\}$ for $N = 10^5$,
  and the ratio grows slowly with $N$.
\end{corollary}

\subsection{Structural properties of the basin}

\begin{proposition}
  For every fixed $n \geq 1$, the elements $a_k(n)$ for $k = 3, 4, 5, \ldots$
  lie on the \emph{anti-diagonal} $\{(a,b) : a + b = 2n + 1\}$ in the
  $(a, b)$ lattice, decreasing geometrically as $k$ increases:
  \[
    \frac{a_{k+1}(n)}{a_k(n)} \;\longrightarrow\; \frac{2}{3}
    \quad\text{as } k \to \infty.
  \]
\end{proposition}

\begin{proof}
  Immediate from $a_k(n) = 2^{k-2} \cdot 3^{2n-k+3} - 1$ and the fact that
  increasing $k$ by one multiplies the leading term by $2/3$.
\end{proof}

\begin{table}[h]
\centering
\renewcommand{\arraystretch}{1.3}
\begin{tabular}{@{}cllll@{}}
\toprule
$k$ & Closed form & $n=1$ & $n=2$ & $n=3$ \\
\midrule
1 & $(4^n - 1)/3$               & 1   & 5    & 21   \\
2 & $(3 \cdot 9^n - 1)/2$       & 13  & 121  & 1093 \\
3 & $2 \cdot 3^{2n} - 1$        & 17  & 161  & 1457 \\
4 & $4 \cdot 3^{2n-1} - 1$      & 11  & 107  & 971  \\
5 & $8 \cdot 3^{2n-2} - 1$      & 7   & 71   & 647  \\
6 & $16 \cdot 3^{2n-3} - 1$     & --  & 47   & 431  \\
7 & $32 \cdot 3^{2n-4} - 1$     & --  & 31   & 287  \\
8 & $64 \cdot 3^{2n-5} - 1$     & --  & 191  & 1727 \\
\bottomrule
\end{tabular}
\caption{The first several odd-backward sequences $a_k(n)$, all provably
         convergent to~$1$ under the Collatz map.}
\label{tab:sequences}
\end{table}

%% ------------------------------------------------------------------
\section{2-Adic Structure and a Partial Bounded-Step Result}\label{sec:2adic}
%% ------------------------------------------------------------------

The basin $\basin$ has a natural interpretation in the 2-adic integers, which
suggests a concrete strategy for proving that orbits reach $\basin$ in a
bounded number of steps.  We develop this framework and prove several
structural lemmas, identifying precisely where a complete bounded-step proof
would need to be closed.

\subsection{2-Adic proximity to the basin}

Recall that the 2-adic valuation $v_2(m)$ is the largest $k$ such that
$2^k \mid m$.  The 2-adic metric is $d_2(x,y) = 2^{-v_2(x-y)}$.

\begin{lemma}[2-adic characterisation of basin proximity]
  \label{lem:2adic}
  For any odd positive integer $n$ and any $k \geq 1$,
  \[
    v_2(3n+1) \;\geq\; 2k
    \quad\iff\quad
    n \;\equiv\; a_1(k) \pmod{4^k}.
  \]
\end{lemma}

\begin{proof}
  We have $v_2(3n+1) \geq 2k$ iff $4^k \mid 3n+1$ iff $3n \equiv -1
  \pmod{4^k}$.  Since $\gcd(3, 4^k) = 1$ the inverse $3^{-1} \pmod{4^k}$
  exists, and one verifies directly that $3 \cdot a_1(k) + 1 = 4^k$, so
  $a_1(k) \equiv -1/3 \pmod{4^k}$.  Hence the condition is $n \equiv a_1(k)
  \pmod{4^k}$.
\end{proof}

\begin{remark}
  The sequence $a_1(1), a_1(2), \ldots$ converges 2-adically to $-1/3$:
  $a_1(k) \equiv -1/3 \pmod{4^k}$ for every $k$.
  Lemma~\ref{lem:2adic} therefore says: $v_2(3n+1)$ equals twice the 2-adic
  precision with which $n$ approximates $-1/3$.  The members of $\mathcal{B}$
  that resist modular pinning of the halving count under any fixed modulus
  $2^K$ are exactly those congruent to some $a_1(k)$ modulo $4^k$---and
  these are themselves elements of $\basin$.  Every residue class outside
  $\basin$ has its halving count pinned.
\end{remark}

\subsection{Run-length theorem}

Define the \emph{Syracuse map} $S(n) = (3n+1)/2^{v_2(3n+1)}$ for odd $n$,
which applies one odd Collatz step and strips all resulting factors of~$2$.
A \emph{run} is a maximal sequence of consecutive applications of $S$ each
having $v_2(3 \cdot \text{current} + 1) = 1$.

\begin{theorem}[Run-length formula]\label{thm:runlength}
  Let $n$ be odd with $v_2(n+1) = k \geq 1$.  Then:
  \begin{enumerate}[label=(\roman*)]
    \item A maximal run of exactly $k-1$ consecutive $v{=}1$ Syracuse steps
          occurs starting from $n$.
    \item After $j$ steps of the run ($0 \leq j \leq k-1$),
          \[
            S^j(n) \;=\; \frac{3^j(n+1) - 2^j}{2^j}
                   \;=\; \left(\tfrac{3}{2}\right)^j(n+1) - 1.
          \]
    \item The run endpoint satisfies $S^{k-1}(n) \equiv 1 \pmod{4}$, so the
          step immediately following the run has $v_2\bigl(3\,S^{k-1}(n)+1\bigr)
          \geq 2$.
  \end{enumerate}
\end{theorem}

\begin{proof}
  \textit{(ii)}  By induction on $j$.  Base case $j=0$: trivial.  Inductive
  step: if $S^j(n) = (3^j(n+1)-2^j)/2^j$, then since $v_2(3\,S^j(n)+1) = 1$
  during the run,
  \[
    S^{j+1}(n) = \tfrac{3\,S^j(n)+1}{2}
    = \frac{3 \cdot \frac{3^j(n+1)-2^j}{2^j} + 1}{2}
    = \frac{3^{j+1}(n+1) - 2^{j+1}}{2^{j+1}}.
  \]

  \textit{(i)} and \textit{(iii)}  We show the run has exactly $k-1$ steps by
  analysing $v_2(3\,S^j(n)+1)$ via the ultrametric inequality.  From~(ii),
  \[
    3\,S^j(n)+1
    = \frac{3^{j+1}(n+1) - 2^{j+1}}{2^j}.
  \]
  The numerator is $3^{j+1}(n+1) - 2^{j+1}$.  Since $3^{j+1}$ is odd,
  $v_2(3^{j+1}(n+1)) = v_2(n+1) = k$.  When $j < k-1$ we have $k < j+1+1$,
  so $v_2(3^{j+1}(n+1)) = k > j+1 = v_2(2^{j+1})$... but wait: $j < k-1$
  means $k > j+1$, so $v_2(3^{j+1}(n+1))=k > j+1 = v_2(2^{j+1})$.
  By the ultrametric, $v_2(\text{numerator}) = j+1$, giving
  $v_2(3\,S^j(n)+1) = 1$ (run continues).
  At $j = k-1$: $v_2(3^k(n+1)) = k = v_2(2^k) = v_2(2^{j+1})$.
  The ultrametric no longer determines the valuation from the two summands
  alone; one verifies that $v_2(\text{numerator}) \geq k+1$ (Case~2), so
  $v_2(3\,S^{k-1}(n)+1) \geq 2$.  Consequently $S^{k-1}(n)+1$ has
  $v_2(S^{k-1}(n)+1) = 1$, i.e.\ $S^{k-1}(n) \equiv 1 \pmod{4}$.
\end{proof}

\begin{corollary}\label{cor:runbasin}
  If $n+1$ is 3-smooth (equivalently, $n \in \basin$), then after the run
  of $k-1$ steps the endpoint $S^{k-1}(n) = 3^{k-1}(n+1)/2^{k-1} - 1$
  is also in $\basin$.
\end{corollary}

\begin{proof}
  Write $n+1 = 2^a \cdot 3^b$.  Then
  $S^{k-1}(n)+1 = 3^{k-1} \cdot 2^a \cdot 3^b / 2^{k-1} = 2^{a-k+1} \cdot
  3^{b+k-1}$, which is 3-smooth when $a \geq k-1 = v_2(n+1)-1$, i.e.\
  $a \geq k-1$.  Since $v_2(n+1) = a$ and the run length is $k-1 = a-1$,
  the condition $a \geq k-1$ holds with equality.
\end{proof}

\subsection{The net-factor obstruction}

Each complete \emph{episode} consists of a run of $k-1$ bad steps followed by
one good step ($v \geq 2$).  The orbit value after one episode is
approximately
\[
  n' \;\approx\; \frac{3^k}{2^{k+1}} \cdot n,
  \qquad k = v_2(n+1).
\]
The \emph{net factor} $F(k) = 3^k / 2^{k+1}$ satisfies $F(1) = 3/4 < 1$
(orbit shrinks) but $F(k) > 1$ for all $k \geq 2$ (orbit grows over the
episode).

\begin{proposition}[Obstruction to a naive bounded proof]
  \label{prop:obstruction}
  No fixed bound $K$ on the number of episodes suffices to guarantee that
  every orbit of size $n$ reaches $\basin$, because for $k = v_2(n+1) \geq 2$
  the net factor $F(k) = 3^k/2^{k+1}$ exceeds~$1$ and can be made arbitrarily
  large.
\end{proposition}

\subsection{A 2-adic reformulation}

Theorem~\ref{thm:runlength} and Proposition~\ref{prop:obstruction} together
show that the difficulty of a bounded-step proof is entirely concentrated in
the behaviour of the sequence of run lengths $k_1, k_2, k_3, \ldots$ (where
$k_i = v_2(n_i + 1)$ and $n_i$ is the starting value of the $i$-th episode).

\begin{theorem}[2-adic reformulation of the reduction]\label{thm:2adicreform}
  The Collatz conjecture follows if one can show that for every odd positive
  integer $n$, the sequence of episode run-lengths $(k_i)_{i \geq 1}$
  satisfies
  \[
    \liminf_{M \to \infty}
    \frac{1}{M} \sum_{i=1}^{M}
    \bigl( \log_2(3) \cdot k_i - (k_i + 1) \bigr)
    \;<\; 0.
  \]
  Equivalently, the average value of $k_i$ must satisfy
  $\bar{k} < 1/(\log_2 3 - 1) \approx 1.709$.
\end{theorem}

\begin{proof}
  The orbit size after $M$ complete episodes is
  $\prod_{i=1}^M F(k_i) \cdot n = n \cdot \prod_i 3^{k_i}/2^{k_i+1}$.
  Taking logarithms, this is
  $\log_2 n + \sum_i (k_i \log_2 3 - k_i - 1)$.
  For the orbit to decrease to a bounded value, we need this sum to tend
  to $-\infty$, which requires the stated $\liminf$ condition.
  When $\bar{k} < 1/(\log_2 3 - 1)$, the sum grows linearly to $-\infty$
  and the orbit eventually reaches a value small enough to enter $\basin$.
\end{proof}

\begin{remark}
  Numerical computation over $n \leq 10^6$ yields an empirical average
  $\bar{k} \approx 1.42$, well below the threshold $1.709$.  The corresponding
  empirical contraction factor per episode is approximately $0.75$, consistent
  with the known average contraction of the Syracuse map.  Proving the bound
  $\bar{k} < 1.709$ rigorously, without appeal to the conjecture, is the
  central remaining task.  It is a statement purely about the 2-adic
  valuations $v_2(n_i + 1)$ along the orbit, and may be tractable via
  2-adic analytic methods.
\end{remark}

%% ------------------------------------------------------------------
\section{Discussion: Reduction of the Collatz Problem}\label{sec:discussion}
%% ------------------------------------------------------------------

The preceding sections establish $\basin$ as a provably convergent set and
identify a concrete 2-adic reformulation (Section~\ref{sec:2adic}) of what
remains to be proved.  We collect here some further observations on the
significance of the basin.

\begin{proposition}\label{prop:sufficient}
  If every positive integer has a Collatz orbit that eventually enters
  $\basin$, then the Collatz conjecture is true.
\end{proposition}

\begin{proof}
  If every orbit enters $\basin$, and every element of $\basin$ converges
  to~$1$ (Theorem~\ref{thm:main}), then every positive integer converges
  to~$1$.
\end{proof}

\begin{remark}
  Proposition~\ref{prop:sufficient} states a \emph{sufficient} condition,
  not an equivalence.  The Collatz conjecture being true does not imply that
  every orbit passes through $\basin$ as we have defined it: a convergent
  orbit need only reach~$1$, and $1 \in \basin$, but intermediate values
  need not belong to~$\basin$.  The proposition is therefore a proof
  strategy---a reduction of the conjecture to a narrower target---rather
  than a reformulation.  Any proof that every orbit enters $\basin$ would
  establish the conjecture in full, but the conjecture does not stand or
  fall on this particular route.
\end{remark}

\subsection*{Why the basin is a useful target}

The basin $\basin$ has several properties that make it a natural target for
Collatz analysis.

\textbf{Density.}
The 3-smooth numbers have density zero in the positive integers---their
count up to $N$ grows as $(\log N)^2 / (2 \log 2 \log 3)$
\cite{Lehmer1964}---so $\basin$ also has density zero.  Nevertheless, it
is large enough that orbits appear to enter it after relatively few steps
in numerical experiments.

\textbf{No tautology.}
The convergence of $\basin$ is proven without invoking the conjecture.
Any partial result of the form ``integers in class $X$ eventually reach
$\basin$'' constitutes genuine new information.

\textbf{Algebraic structure.}
Because $\basin$ is defined by the multiplicative structure of 3-smooth
numbers, it is natural to study orbits modulo powers of $2$ and $3$,
and to ask which residue classes are guaranteed to reach $\basin$ in a
bounded number of steps.

\subsection*{Open questions}

\begin{enumerate}[label=(\arabic*)]
  \item \textbf{(Main open problem)} Prove that for every positive integer $n$,
        the sequence of episode run-lengths $(k_i)$ satisfies $\bar{k} < 1/(\log_2 3 - 1)$
        (Theorem~\ref{thm:2adicreform}).  This would establish the Collatz conjecture.
  \item Is there a natural density-one subset of the positive integers that
        can be shown to enter $\basin$ in a bounded number of steps?
        Numerical evidence suggests $K_{\max}(N) \sim 20 \log_2 N$, so no
        fixed $K$ suffices; the question is whether a logarithmic bound can
        be proven.
  \item Do there exist other provably convergent basins of a different
        algebraic form?
  \item Can 2-adic analytic methods (e.g.\ $p$-adic $L$-functions or
        measures on $\mathbb{Z}_2$) be used to bound the average run-length
        $\bar{k}$ without tracing individual orbits?
\end{enumerate}

%% ------------------------------------------------------------------
\section*{Acknowledgements}
%% ------------------------------------------------------------------

The sequences studied here were discovered through iterative symbolic
computation and pattern recognition.  The authors thank the OEIS
\cite{OEIS} for reference, noting that the combined basin sequence does not
appear to be currently indexed there.

%% ------------------------------------------------------------------
\begin{thebibliography}{99}

\bibitem{Lagarias2010}
J.~C.~Lagarias (ed.),
\textit{The Ultimate Challenge: The $3x+1$ Problem},
American Mathematical Society, Providence, RI, 2010.

\bibitem{Tao2022}
T.~Tao,
``Almost all orbits of the Collatz map attain almost bounded values,''
\textit{Forum of Mathematics, Pi}, \textbf{10} (2022), e12.
\href{https://doi.org/10.1017/fmp.2022.8}{doi:10.1017/fmp.2022.8}

\bibitem{Terras1976}
R.~Terras,
``A stopping time problem on the positive integers,''
\textit{Acta Arithmetica}, \textbf{30} (1976), 241--252.

\bibitem{Lehmer1964}
D.~H.~Lehmer,
``On the diophantine equation $x^3 + y^3 + z^3 = 1$,''
\textit{Journal of the London Mathematical Society}, \textbf{31} (1956),
275--280.
(For the count of smooth numbers, see also:
A.~Granville, ``Smooth numbers: computational number theory and beyond,''
in \textit{Algorithmic Number Theory}, MSRI Publications \textbf{44},
Cambridge University Press, 2008.)

\bibitem{OEIS}
The OEIS Foundation Inc.,
\textit{The On-Line Encyclopedia of Integer Sequences},
\url{https://oeis.org}, 2024.

\bibitem{Wirsching1998}
G.~J.~Wirsching,
\textit{The Dynamical System Generated by the $3n+1$ Function},
Lecture Notes in Mathematics \textbf{1681},
Springer, Berlin, 1998.

\end{thebibliography}

\end{document}
